%
% teil6.tex -- Beweis des Satzes von Napoleon
%
\section{Beweis des Satzes von Napoleon
\label{komplex:section:beweis}}
\kopfrechts{Beweis des Satzes von Napoleon}
In diesem Abschnitt wird der Satz von Napoleon präzise formuliert
und mit den bisher entwickelten Werkzeugen vollständig bewiesen.
Die Unterabschnitte gliedern den Beweis in Formulierung, Beweisidee,
algebraische Durchführung und abschliessende Diskussion.

%
% Formulierung des Satzes
%
\subsection{Formulierung des Satzes
\label{komplex:subsection:formulierung}}
\begin{satz}[Napoleon]
\label{komplex:satz:napoleon}
\index{Satz!von Napoleon}%
Sei $ABC$ ein Dreieck in der Ebene, dessen Eckpunkte als komplexe Zahlen
$A, B, C$ mit positivem Drehsinn nummeriert sind.
Werden über jeder Seite gleichseitige Dreiecke nach aussen errichtet
und deren Schwerpunkte mit $O_a, O_b, O_c$ bezeichnet,
so ist das Dreieck $O_a O_b O_c$ gleichseitig.
\end{satz}

%
% Beweisidee und Charakterisierung der Gleichseitigkeit
%
\subsection{Beweisidee und Charakterisierung der Gleichseitigkeit
\label{komplex:subsection:beweisidee}}
Die Beweisstrategie nutzt eine besonders kompakte Charakterisierung
gleichseitiger Dreiecke mittels des Drehfaktors $d$ aus
\eqref{komplex:eq:d60}:
Ein Dreieck mit Eckpunkten $u, v, w$ ist genau dann gleichseitig,
wenn man den Vektor von $w$ nach $v$ durch Drehung des Vektors von $w$
nach $u$ um $60^\circ$ erhält, also wenn
\begin{equation}
(v - w)
=
d \cdot (u - w) ,
\qquad \text{bzw.} \qquad
(v - w) - d \cdot (u - w) = 0 .
\label{komplex:eq:charakterisierung}
\end{equation}
Diese Bedingung kodiert beides zugleich:
Da $|d| = 1$, sind die Strecken $|v - w|$ und $|u - w|$ gleich lang;
und da $\arg(d) = 60^\circ$, schliessen sie einen Winkel von $60^\circ$
ein --- der dritte Innenwinkel ist damit ebenfalls $60^\circ$,
das Dreieck also gleichseitig.

Der Beweis besteht im Wesentlichen darin,
die drei Schwerpunkte aus \eqref{komplex:eq:schwerpunkte} in den Ausdruck
\begin{equation}
\delta
=
(O_b - O_c) - d \cdot (O_a - O_c)
\label{komplex:eq:delta}
\end{equation}
einzusetzen und mit Hilfe der Identität
$d^2 - d + 1 = 0$ aus \eqref{komplex:eq:d-identitaet} zu zeigen,
dass $\delta = 0$ ist.

%
% Algebraische Durchfuehrung
%
\subsection{Algebraische Durchführung
\label{komplex:subsection:algebraische-durchfuehrung}}
Wir schreiben zunächst die drei Schwerpunkte explizit als Linearkombinationen
in $A, B, C$.
Mit den Verbindungsstrecken aus \eqref{komplex:eq:verbindungsstrecken} und
den Definitionen \eqref{komplex:eq:p-punkte} der Punkte $P_a, P_b, P_c$
ergibt sich
\begin{align}
O_a
&=
\frac{B + C + P_a}{3}
=
\frac{B + C + C + d (B - C)}{3}
=
\frac{(1 + d) B + (2 - d) C}{3} ,
\notag \\
O_b
&=
\frac{A + C + P_b}{3}
=
\frac{A + C + A + d (C - A)}{3}
=
\frac{(2 - d) A + (1 + d) C}{3} ,
\label{komplex:eq:o-explizit}
\\
O_c
&=
\frac{A + B + P_c}{3}
=
\frac{A + B + B + d (A - B)}{3}
=
\frac{(1 + d) A + (2 - d) B}{3} .
\notag
\end{align}
Aus diesen Ausdrücken berechnen wir die beiden Differenzen
\begin{align}
O_b - O_c
&=
\frac{(2 - d) A + (1 + d) C - (1 + d) A - (2 - d) B}{3}
\notag \\
&=
\frac{(1 - 2d) A - (2 - d) B + (1 + d) C}{3} ,
\label{komplex:eq:Ob-Oc}
\\[1ex]
O_a - O_c
&=
\frac{(1 + d) B + (2 - d) C - (1 + d) A - (2 - d) B}{3}
\notag \\
&=
\frac{-(1 + d) A + (2d - 1) B + (2 - d) C}{3} .
\label{komplex:eq:Oa-Oc}
\end{align}
Damit lässt sich der Ausdruck $\delta$ aus \eqref{komplex:eq:delta}
nach $A, B, C$ sortieren:
\begin{equation}
\delta
=
\frac{1}{3}
\Bigl(
( 1 - 2d + d (1 + d) ) A
- ( 2 - d + d (2d - 1) ) B
+ ( 1 + d - d (2 - d) ) C
\Bigr) .
\label{komplex:eq:delta-sortiert}
\end{equation}
Die drei Klammerausdrücke ergeben genau denselben quadratischen Ausdruck in $d$:
\begin{align*}
1 - 2d + d (1 + d)
&=
1 - 2d + d + d^2
=
d^2 - d + 1 ,
\\
2 - d + d (2d - 1)
&=
2 - d + 2 d^2 - d
=
2 (d^2 - d + 1) ,
\\
1 + d - d (2 - d)
&=
1 + d - 2 d + d^2
=
d^2 - d + 1 .
\end{align*}
Setzen wir diese drei Ergebnisse in \eqref{komplex:eq:delta-sortiert} ein,
so erhalten wir
\begin{equation}
\delta
=
\frac{(A - 2 B + C) \cdot (d^2 - d + 1)}{3} .
\label{komplex:eq:delta-faktorisiert}
\end{equation}
Nach der Identität \eqref{komplex:eq:d-identitaet} ist aber
$d^2 - d + 1 = 0$, und damit folgt $\delta = 0$,
unabhängig von der konkreten Wahl von $A, B, C$.
Die Bedingung \eqref{komplex:eq:charakterisierung} ist erfüllt,
und das Dreieck $O_a O_b O_c$ ist gleichseitig.
\qed

%
% Diskussion der Beweisstrategie
%
\subsection{Diskussion der Beweisstrategie
\label{komplex:subsection:diskussion}}
Der Beweis ist bemerkenswert kurz, wenn man die zugrunde liegenden
Identitäten kennt.
Sein Erfolg beruht auf drei zusammenwirkenden Beobachtungen:
erstens auf der direkten algebraischen Charakterisierung gleichseitiger
Dreiecke durch die Drehbedingung \eqref{komplex:eq:charakterisierung},
zweitens auf der Tatsache, dass die Drehung um $60^\circ$ als
Multiplikation mit einer einzigen komplexen Zahl darstellbar ist,
und drittens auf der Identität $d^2 - d + 1 = 0$, die alle drei
Koeffizienten in \eqref{komplex:eq:delta-sortiert} simultan zum
Verschwinden bringt.

Im Vergleich dazu ist ein rein synthetisch-geometrischer Beweis deutlich
aufwändiger und benötigt mehrere Hilfskonstruktionen,
wie sie etwa in~\cite{komplex:coxeter} dargestellt werden.
Die algebraische Methode tauscht geometrische Anschauung gegen
rechnerische Eleganz --- ein Tausch, der im konkreten Fall zugunsten
der Algebra ausfällt.

Bemerkenswert ist auch, dass der Beweis vollständig linear in den
Eckpunkten $A, B, C$ ist.
Dies bedeutet, dass das Resultat unter Ähnlichkeitstransformationen ---
\index{Ahnlichkeitstransformation@Ähnlichkeitstransformation}%
Translation, Drehung und Streckung --- invariant bleibt,
eine Eigenschaft, die mit synthetischen Methoden nicht so unmittelbar
sichtbar wird.
Die kompakte Faktorisierung \eqref{komplex:eq:delta-faktorisiert} zeigt
zudem, dass die Gleichseitigkeit allein aus der Geometrie der sechsten
Einheitswurzel folgt, also aus einer einzigen algebraischen Tatsache über
\index{Einheitswurzel}%
$d$.
